\scanpage{25}
\usetikzlibrary{shapes.geometric}
% archive sheet number: 24
% Whole draft cancelled by two long diagonal strokes.
\margin{\mbox{\tikz[baseline=(c.base)]{\node[draw,ellipse,inner sep=1pt] (c) {\deleted{\unsure{$G$ comme dessus}}};}}}

% Source lines 1--7.
\deleted{\uline{Théorème} Soit $H$ un sous-grpe de $G$, simple/$S$ et à fibres connexes, ayant $\unsure{=|\text{même}}$ rang réductif et $=$ \unsure{rg|rang} \unsure{unipotent|unip.} que $G$, \deleted{Alors} \deleted{\unsure{$H$ est le}} \unsure{soit} $\unsure{\mathfrak h|\mathscr H}=\operatorname{Lie}(H)\subset\unsure{\mathfrak g|\mathscr G}=\operatorname{Lie}(G)$. \deleted{Alors} Alors $H$ est le \uline{normalisateur connexe} de $\mathfrak h$ dans $G$ \deleted{\unsure{(qui a le rg}} (opérant sur $\mathfrak g$ par la représentation adjointe).}

% Source lines 8--16.
\deleted{En effet, Demazure a montré qu'il suffit de prouver que pour tt fibre géométrique, le normalisateur de $\mathfrak h$ dans \uuline{$\mathfrak g$} est égal à $\mathfrak h$. On peut donc supposer $S$ spectre d'un corps alg. clos, et \deleted{\unsure{on a}} \unsure{notre|not.} hypothèse signifie alors que $H$ contient un sous-groupe de Cartan $C$ \uline{de $G$} ; soit $T$ \unsure{son|un} tore maximal.}
